% Lines 2077-2100 and 5466-5613 from original attention_transformers_notes_ghost.tex
% Appendix: Future Directions

%==============================================================================
\section{Future Directions}
%==============================================================================

\subsection{Research Trends}

\begin{enumerate}
    \item \textbf{Efficiency}: Linear attention, sparse patterns, compression
    \item \textbf{Long Context}: Handling sequences with millions of tokens
    \item \textbf{Multimodal}: Unified architectures for vision, text, and audio
    \item \textbf{Reasoning}: Chain-of-thought, multi-step reasoning
    \item \textbf{Continual Learning}: Adapting without forgetting
\end{enumerate}

\subsection{Open Challenges}

\begin{itemize}
    \item \textbf{Theoretical Understanding}: Why do transformers work so well?
    \item \textbf{Inductive Biases}: What biases should we build in?
    \item \textbf{Sample Efficiency}: Reducing data requirements
    \item \textbf{Interpretability}: Understanding model decisions
    \item \textbf{Robustness}: Handling distribution shift and adversarial inputs
\end{itemize}

% --- Continuation from line 5466 ---

%==============================================================================
\section{Future Directions and Open Problems}
%==============================================================================

\subsection{Theoretical Challenges}

\subsubsection{Understanding Emergent Abilities}

As transformers scale, they exhibit emergent abilities not present in smaller models. Key open questions:

\begin{enumerate}
    \item \textbf{Phase Transitions:} At what scale do specific abilities emerge?
    \item \textbf{Compositional Generalization:} How do transformers learn to compose known concepts?
    \item \textbf{In-Context Learning:} What mechanisms enable few-shot learning without parameter updates?
\end{enumerate}

\subsubsection{Attention Mechanistic Interpretability}

Understanding what attention heads learn:

\begin{lstlisting}[language=Python, caption=Attention Head Analysis]
class AttentionHeadAnalyzer:
    """Analyze specialization of attention heads"""

    def __init__(self, model: nn.Module) -> None:
        self.model = model
        self.head_patterns: dict[str, list[Tensor]] = {
            'positional': [],
            'syntactic': [],
            'semantic': []
        }

    def analyze_positional_heads(
        self,
        attention_weights: Tensor
    ) -> Tensor:
        """Identify heads attending to specific positions"""

        batch_size, n_heads, seq_len, seq_len = attention_weights.shape

        # Compute average attention pattern
        avg_pattern = attention_weights.mean(dim=0)

        # Analyze diagonal patterns (attending to same position)
        diagonal_score = torch.zeros(n_heads)
        for h in range(n_heads):
            diagonal = torch.diagonal(avg_pattern[h])
            diagonal_score[h] = diagonal.mean()

        # Analyze attending to previous token
        prev_token_score = torch.zeros(n_heads)
        for h in range(n_heads):
            prev_diagonal = torch.diagonal(avg_pattern[h], offset=-1)
            prev_token_score[h] = prev_diagonal.mean()

        # Analyze attending to first token
        first_token_score = avg_pattern[:, :, 0].mean(dim=1)

        return torch.stack([
            diagonal_score,
            prev_token_score,
            first_token_score
        ], dim=1)

    def analyze_syntactic_heads(
        self,
        attention_weights: Tensor,
        dependency_matrix: Tensor
    ) -> Tensor:
        """Identify heads capturing syntactic dependencies"""

        # Compare attention patterns with syntactic dependencies
        batch_size, n_heads, seq_len, seq_len = attention_weights.shape

        syntactic_scores = torch.zeros(n_heads)

        for h in range(n_heads):
            # Compute correlation with dependency matrix
            attn_flat = attention_weights[:, h].flatten()
            dep_flat = dependency_matrix.flatten()

            # Pearson correlation
            correlation = torch.corrcoef(
                torch.stack([attn_flat, dep_flat])
            )[0, 1]

            syntactic_scores[h] = correlation

        return syntactic_scores
\end{lstlisting}

\subsection{Architectural Innovations}

\subsubsection{Continuous Attention}

Replace discrete attention with continuous formulations:

\begin{equation}
\text{ContinuousAttention}(q, K, V) = \int_{0}^{1} \alpha(q, k(t)) \cdot v(t) \, dt
\end{equation}

where $k(t)$ and $v(t)$ are continuous functions of position.

\subsubsection{Quantum-Inspired Attention}

Leverage quantum computing concepts:

\begin{lstlisting}[language=Python, caption=Quantum-Inspired Attention Concept]
class QuantumInspiredAttention(nn.Module):
    """Attention using quantum superposition concepts"""

    def __init__(self, d_model: int, n_qubits: int) -> None:
        super().__init__()
        self.d_model = d_model
        self.n_qubits = n_qubits
        self.n_states = 2 ** n_qubits

        # Quantum state embeddings
        self.state_embed = nn.Parameter(
            torch.randn(self.n_states, d_model)
        )

        # Measurement operators
        self.measurement = nn.Linear(d_model, self.n_states)

    def forward(
        self,
        query: Tensor,
        key: Tensor,
        value: Tensor
    ) -> Tensor:
        # Create quantum superposition
        q_probs = F.softmax(self.measurement(query), dim=-1)
        k_probs = F.softmax(self.measurement(key), dim=-1)

        # Quantum interference pattern
        interference = torch.einsum('bih,bjh->bij', q_probs, k_probs)

        # Collapse to classical attention
        attention_weights = interference.sum(dim=-1, keepdim=True)
        attention_weights = F.softmax(attention_weights, dim=-2)

        # Apply attention
        output = torch.matmul(attention_weights, value)

        return output
\end{lstlisting}
